\chapter{Sieves as probability}\label{ch:oc-sieves}

This chapter reads a sieve as an inclusion--exclusion expectation, and isolates the parity barrier that stops the method from detecting primes.

\section{Inclusion--exclusion as expectation}

The sieve weights that count integers free of small prime factors are partial sums of $\prod_{p<z}(1-\mathbf 1_{p\mid n})$. Truncating the product is the same as truncating an expectation.

\section{The parity barrier}

Selberg's parity phenomenon says that a sieve of this classical type cannot distinguish integers with an even number of prime factors from integers with an odd number. In particular it cannot, by itself, prove that a polynomial takes infinitely many prime values.
